\documentclass[12pt]{article}
\usepackage[margin=0.6in,top=0.85in,bottom=0.55in]{geometry}
\usepackage{lmodern}
\usepackage{microtype}
\usepackage{fancyhdr}
\usepackage{titlesec}
\usepackage{enumitem}
\usepackage{lastpage}
\usepackage{array}
\usepackage[hidelinks]{hyperref}

\setlength{\parindent}{0pt}
\setlength{\parskip}{0.25em}
\setlength{\headheight}{26pt}

\titleformat{\section}{\large\bfseries\scshape}{}{0pt}{}
\titleformat{\subsection}{\normalsize\bfseries}{}{0pt}{}

\pagestyle{fancy}
\fancyhf{}
\fancyhead[C]{%
\footnotesize
Student Name: \hspace{2.3in} \quad Assignment ID: U1L1AW \quad Page \thepage\ of \pageref{LastPage}
}
\renewcommand{\headrulewidth}{0.5pt}
\renewcommand{\footrulewidth}{0pt}

\newcommand{\answerbox}[2]{% #1 = label e.g. "Part 1 Q1 ANSWER BOX", #2 = height e.g. 2.6in
\vspace{0.15em}
\noindent\textbf{#1}\par
\vspace{0.05em}
\noindent\fbox{%
\begin{minipage}[t][#2]{\dimexpr\textwidth-2\fboxsep-2\fboxrule}
\rule{\textwidth}{0pt}
\vspace{#2}
\end{minipage}
}\par
\vspace{0.25em}
}

\begin{document}

\begin{center}
{\LARGE\bfseries Macbeth Anchor Worksheet}\par\vspace{0.05em}
{\large\scshape Week 1: Prophecy Is Not Proof}\par\vspace{0.1em}
\rule{0.78\textwidth}{0.5pt}\par\vspace{0.15em}
\end{center}

\noindent\textbf{Purpose:} This worksheet asks you to use the paired \textit{Macbeth Lit-Anchor Reading} to test Macbeth's reasoning with Whately's four questions: conclusion, stated reason, hidden assumption, and whether the conclusion follows. Use evidence from both Macbeth passages. Your answers should analyze reasoning, not simply retell the plot.

\vspace{0.3em}

\section*{Part 1: Prophecy and Inference}

\textbf{Part 1 Q1:} In Act 1, Scene 3, what exactly do the witches tell Macbeth, and what do they \textbf{not} tell him? Explain the difference between a prediction and a command.

\answerbox{Part 1 Q1 ANSWER BOX}{1.75in}

\textbf{Part 1 Q2:} Macbeth treats the two confirmed titles as ``happy prologues'' to kingship. Use Whately's four-part test to analyze Macbeth's reasoning: identify the conclusion, stated reason, hidden assumption, and whether the conclusion follows.

\answerbox{Part 1 Q2 ANSWER BOX}{1.85in}

\newpage

\section*{Part 2: Banquo's Warning}

\textbf{Part 2 Q1:} Why is Banquo's question, ``What, can the devil speak true?'' logically important? Explain how Banquo gives Macbeth a rival interpretation of the prophecy.

\answerbox{Part 2 Q1 ANSWER BOX}{2.2in}

\textbf{Part 2 Q2:} Explain this Week 1 claim in your own words: \textbf{prediction is not permission; possibility is not command; desire is not proof}. Use one moment from Act 1, Scene 3 as evidence.

\answerbox{Part 2 Q2 ANSWER BOX}{2.25in}

\newpage

\section*{Part 3: Macbeth Weighs Reasons}

\textbf{Part 3 Q1:} In Act 1, Scene 7, what conclusion does Macbeth's soliloquy seem to support? List at least three reasons Macbeth gives against murdering Duncan.

\answerbox{Part 3 Q1 ANSWER BOX}{2.35in}

\textbf{Part 3 Q2:} Macbeth says he has ``no spur'' except ``vaulting ambition.'' Is ambition a reason, a desire, or both? Explain why ambition alone does or does not prove that Macbeth should act.

\answerbox{Part 3 Q2 ANSWER BOX}{2.0in}

\newpage

\section*{Part 3: Macbeth Weighs Reasons, Continued}

\textbf{Part 3 Q3:} Test this argument: ``Macbeth should not murder Duncan.'' Use Macbeth's reasons from Act 1, Scene 7 and explain whether the conclusion follows.

\answerbox{Part 3 Q3 ANSWER BOX}{2.2in}

\section*{Part 4: Anchor Synthesis}

\textbf{Part 4 Q1:} Choose one claim below, then complete a four-part Macbeth argument skeleton. Your answer must include \textbf{Claim/conclusion}, \textbf{Stated reason}, \textbf{Hidden assumption}, and \textbf{Follows}.\par
\vspace{0.1em}
\begin{itemize}[leftmargin=*, itemsep=0.08em, topsep=0.08em, parsep=0pt]
\item The prophecy proves Macbeth will be king.
\item If Macbeth is destined to be king, he is not responsible for what he does.
\item Macbeth should not kill Duncan.
\item Ambition is Macbeth's only real reason for murder.
\end{itemize}

\answerbox{Part 4 Q1 ANSWER BOX}{2.0in}

\newpage

\section*{Part 4: Anchor Synthesis, Continued}

\textbf{Part 4 Q2:} Write one strong paragraph answering this question: What does Macbeth show about Whately's claim that logic can discipline reasoning but cannot replace virtue? Use at least one specific moment from the reading.

\answerbox{Part 4 Q2 ANSWER BOX}{3.0in}

\end{document}
